% SPDX-License-Identifier: CC-BY-SA-4.0
% Author: Matthieu Perrin
% Part: <Nom de la partie>
% Section: <Nom de la section>
% Sub-section: <Nom de la sous-section>  % (facultatif, laisser vide si non utilisé)
% Frame: <Titre de la slide>

\begingroup

\begin{frame}{Totalité d'un langage algébrique}

  \on[text, top=-3mm]{
    \Probleme{Totality-Alg}{
      Une grammaire algébrique $G$ sur un alphabet $\Sigma$
    }{
      Est-ce que $\mathcal{L}(G) = \Sigma^\star$ ?
    }
  }

  \onBlock[y=10mm]{Théorème (admis)}{
    \textsc{Totality-Alg} est indécidable.
  }

  \onExampleBlock[y=-8mm]{Exemples}{
    \begin{itemize}
    \item Instance positive : \example{$\left\{ S \rightarrow a S \mid b  S \mid \varepsilon\right.$}
    \item Instance négative : \example{$\left\{ S \rightarrow a S b \mid \varepsilon\right.$}
    \end{itemize}
  }

  \onAlertBlock[bottom=3mm]{Remarque -- restriction aux langages déterministes}{
    \small
    \begin{itemize}
    \item \textsc{Totality-Alg} est \alert{indécidable} pour les automates à pile \alert{non-déterministes}
    \item \textsc{Totality-Alg} est \alert{décidable} pour les automates à pile \alert{déterministes}
    \end{itemize}
  }
  
  \footnoteref{S. Ginsburg \& H. G. Rice. \emph{Two Families of Languages Related to ALGOL}. JACM (1962)}
  
\end{frame}

\endgroup
